% !TeX program = xelatex
% UTF-8. Compile with XeLaTeX from this directory or the repository root.
% Source: Google Drive / Notebook / neonate evaluation, in filename order.
% pic01.jpg = 20261018_193408555.jpg; Drive ID 1ppTI44wsK54Txf0z9VU_KJIcb8RHJdQU.
% pic02.jpg = 20261018_193412869.jpg; Drive ID 1MNgHlY7kGqpNd-tOX6Jf689wfmGSX-Dg.
% Define \NeonateEvaluationInNotebook before inputting into a XePersian document.
\ifdefined\NeonateEvaluationInNotebook
  \let\NEfinish\relax
\else
  \def\NEfinish{\end{document}}
  \documentclass[12pt,a4paper]{article}
  \usepackage[margin=22mm]{geometry}
  \usepackage{amsmath,amssymb}
  \usepackage{graphicx}
  \usepackage{array,tabularx}
  \usepackage{enumitem}
  \usepackage{xepersian}
  \IfFontExistsTF{Amiri}{\settextfont{Amiri}}{\settextfont{DejaVu Sans}}
  \setlatintextfont{DejaVu Serif}
  \setlength{\parindent}{0pt}
  \setlength{\parskip}{5pt}
  \setlength{\emergencystretch}{3em}
  \begin{document}
\fi
\providecommand{\NEmark}[1]{\textbf{#1}}

\begin{center}
  {\Large\bfseries ۵۸ ـ ارزیابی نوزاد}
\end{center}

\textbf{یادداشت رونویسی:}
متن و جدول به ترتیب دو تصویر آمده‌اند؛ عبارت‌های مشخص‌شده با ماژیک پررنگ‌اند.
علامت‌های ضربدر جدول مطابق دست‌نویس حفظ شده‌اند.
تصاویر کامل منبع در پایان همین سند آمده‌اند.

\section*{تصویر ۱ ـ صفحهٔ ۱}
\subsection*{آپگار}
\begingroup
\renewcommand{\arraystretch}{1.65}
\setlength{\tabcolsep}{5pt}
\begin{tabularx}{\linewidth}{|>{\raggedleft\arraybackslash}p{0.21\linewidth}|>{\centering\arraybackslash}p{0.09\linewidth}|>{\centering\arraybackslash}X|>{\centering\arraybackslash}X|}
  \hline
  & \textbf{۰} & \textbf{۱} & \textbf{۲} \\
  \hline
  \lr{HR} & ۰ & \lr{\(<100\)} & \lr{\(>100\)} \\
  \hline
  تنفس & \(\times\) & گریهٔ ضعیف & گریهٔ شدید \\
  \hline
  تون عضلانی & \(\times\) & فلکشن جزئی & دست و پا فلکس \\
  \hline
  رفلکس ایریتیبلتی & \(\times\) & حرکات جزئی & گریه، ویتراوال \\
  \hline
  رنگ بدن & آبی & بدن صورتی، اندام آبی & همهٔ بدن صورتی \\
  \hline
\end{tabularx}
\endgroup

\subsection*{اقدامات اتاق زایمان}
\begin{itemize}
  \item پماد نقره نیترات ۱\% برای جلوگیری از کنژکتیویت گنوکوکی.
  \item در مقابل کلامیدیا مؤثر نیست.
\end{itemize}
\textbf{حاشیهٔ متصل به دو سطر بالا:}
«برای هر دو قطره اریترومایسین».

\begin{itemize}
  \item ویتامین \lr{K} عضلانی، پروفیلاکسی برای همه.
  \item واکسن هپاتیت \lr{B} برای همه.
  \item آپگار \NEmark{دقیقهٔ ۱ و ۵} اندازه‌گیری شود.
  \item آپگار ۴ تا ۷: \lr{close attention}.
  \item آپگار ۰ تا ۳: ایست قلبی + تنفسی، برادی‌کاردی شدید.
  \item اگر نوزاد آسفیکسی بگیرد:
        \NEmark{\lr{\(pH<7\)}، تشنج، کما، هایپوتونی}.
  \item ۸۵\% نوزادان خودبه‌خود ظرف ۱۰ تا ۳۰ ثانیه شروع به نفس کشیدن می‌کنند.
  \item ۱۰\% دیگر با کمک شروع به نفس کشیدن می‌کنند.
\end{itemize}

\subsection*{احیای نوزاد}
\begin{itemize}
  \item \NEmark{گرم و خشک کردن نوزاد، اطمینان از باز بودن راه هوایی،
        ساکشن دهان و بینی، صاف کردن نوزاد، تحریک کردن نوزاد}.
  \item اگر به نوزاد بخواهیم اکسیژن بدهیم، باید ۲۱\% باشد.
  \item اگر مؤثر نبود: \lr{\(FiO_2\) to 100}.
  \item با ریت \NEmark{۴۰ تا ۶۰ تنفس در دقیقه}.
\end{itemize}

\begingroup\footnotesize
\textbf{نوشته‌های سربرگ و پابرگ تصویر ۱:}
\lr{Date:}، \lr{Subject:}؛ موضوع: «۵۸ ـ ارزیابی نوزاد»؛ تاریخ خالی است.
نام چاپی \lr{GNOTE} و شمارهٔ دست‌نویس ۱ در پایین برگه.
\par\endgroup

\clearpage
\section*{تصویر ۲ ـ صفحهٔ ۲}
\subsection*{ادامهٔ احیای نوزاد}
\begin{itemize}
  \item اگر برادی‌کاردی به دنبال تنفس مصنوعی بهبود نیافت، یا نبض محیطی لمس نمی‌شود:
        اقدام بعدی ماساژ قلبی است؛ \NEmark{ریت ۱۲۰}،
        نسبت ماساژ به تنفس \lr{\(\frac{3}{1}\)}.
  \item اگر برادی‌کاردی به ونتیلیشن جواب نداد:
        \NEmark{اپی‌نفرین} هم باید بدهیم (\lr{IV})؛
        \NEmark{\lr{0.1--0.3 ml/kg}}، \lr{1:10,000}.
\end{itemize}

\subsection*{\lr{\(SpO_2\)} پس از تولد}
\begin{itemize}
  \item دقیقهٔ ۱: \NEmark{۶۰ تا ۶۵\%}.
  \item دقیقهٔ ۱۰: \NEmark{۸۵ تا ۹۵\%}.
\end{itemize}

\subsection*{علائم ویتراوال نوزاد}
\begin{itemize}
  \item مصرف هروئین مادر: روز ۱ تا ۵ بعد از تولد.
  \item مصرف متادون مادر: هفتهٔ ۱ تا ۴ پس از تولد.
  \item \lr{Sneezing}، اسهال، تب، تهوع، تاکی‌پنه، ترمور، جیترنس،
        گریه‌های جیغ، تشنج، اختلال خواب، \lr{poor feeding}.
  \item \NEmark{اگر مادر متادون مصرف می‌کند: به کودک نالوکسان نمی‌دهیم؛ تشنج}.
\end{itemize}

\begingroup\footnotesize
\textbf{نوشته‌های سربرگ و پابرگ تصویر ۲:}
«۵۸ ـ ارزیابی نوزاد»؛ نام چاپی \lr{GNOTE} و شمارهٔ دست‌نویس ۲ در پایین برگه.
\par\endgroup

\clearpage
\section*{تصاویر کامل منبع}
\IfFileExists{pediatrics/ped-neonatology/neonate evaluation/images/pic01.jpg}{%
  \def\NEimagebase{pediatrics/ped-neonatology/neonate evaluation/images/}}{%
  \def\NEimagebase{images/}}
\begin{center}
  \textbf{تصویر ۱ ـ \lr{pic01.jpg}}\par
  \includegraphics[width=\linewidth,height=0.87\textheight,keepaspectratio]{\NEimagebase pic01.jpg}
\end{center}
\clearpage
\begin{center}
  \textbf{تصویر ۲ ـ \lr{pic02.jpg}}\par
  \includegraphics[width=\linewidth,height=0.87\textheight,keepaspectratio]{\NEimagebase pic02.jpg}
\end{center}
\NEfinish
